\section{The cap-free one-active B/F0 killed-resolvent
theorem}

\textbf{Proof-first candidate, 2026-08-12 PDT.} This note removes the
finite inactive-cap stop from the one-active Bellman/Flat0 argument. It
proves a literal all-clock theorem on the countable inactive phase. No
orientation, reaction word, population box, or stochastic history is
enumerated.

The result is classwise. If the killed inactive phase has a closed
no-kill component, that component is proved positive recurrent directly.
Otherwise its exact killed resolvent is used, and the first access or
Bellman-linkage firing is completed before the episode stops.

\subsection{1. Statement}

Fix two disjoint binary linkage supports \(L_B,L_0\) on species
\(X,U,V\), arbitrary strongly connected orientations, arbitrary positive
labelled rates, and a closed irreducible population class \(\Gamma\).
Assume that \(L_0\) is \textbf{Flat0} in \(X\):

\[
                         y_X=0\qquad(y\in L_0).       \tag{1.1}
\]

Assume that \(L_B\) has a Bellman witness \(q,c\in L_B\):

\[
 q_X=1,\qquad c_X=0,\qquad
 q_U\le c_U,\qquad q_V\le c_V.                       \tag{1.2}
\]

Let \((x_n,t_n)\) be an escaping sequence of physical marked states in
\(\Gamma\), with \(t_n\in L_0\), for which \(X\) is the unique active
coordinate. Thus

\[
 X_n\longrightarrow\infty,\qquad
 \log(1+U_n)+\log(1+V_n)=o(\log X_n),\qquad x_n\ge t_n.
                                                               \tag{1.3}
\]

Carry the actual reaction target as the mark and use

\[
 F(x,t)=\sum_{i\in\{X,U,V\}}\log((x_i-t_i)!),
 \qquad W(x,t)=1+F(x,t).                              \tag{1.4}
\]

\begin{quote}
\textbf{Theorem 1.1 (cap-free B/F0 completion).} Under (1.1)--(1.3), one
of the following two classwise alternatives holds.

\begin{enumerate}
\def\labelenumi{\arabic{enumi}.}
\tightlist
\item
  The full physical class \(\Gamma\) is a fixed-\(X\), fixed-\(U\)
  one-species binary class and is positive recurrent (a singleton is
  allowed).
\item
  A finite support-dependent menu selects, at every sufficiently large
  start in (1.3), an almost surely finite all-clock stopping time
  \(\tau_n>0\), with its actual population endpoint and actual target,
  such that, for constants \(\eta>0\) and \(C_m<\infty\), \[
  \mathbb E_{x_n,t_n}
     [W(X_{\tau_n},T_{\tau_n})-W(x_n,t_n)+\eta\tau_n]
        \longrightarrow-\infty,                    \tag{1.5}
  \] \[
  \mathbb E_{x_n,t_n}
     [((F(X_{\tau_n},T_{\tau_n})-F(x_n,t_n))^+)^m]
        \le C_m.                                    \tag{1.6}
  \] The killed Flat0 prefix has a uniform exponential relative marked
  factorial bound, and its physical duration satisfies the sharp bound
  \[
     \mathbb E\sigma^m\le
        C_m\{1+\log(2+V_n)\}^m.                    \tag{1.7}
  \]
\end{enumerate}
\end{quote}

The rule never stops merely because a cap, tier, active-set label, or
enabled-source flag changes. A Flat0 access jump and a Bellman-linkage
launch are terminal events only after their physical reward completion
has been included.

The logarithm in (1.7) cannot be replaced by a constant uniform in the
inactive start. For the pure-death phase \(V\to0\), followed by a launch
available only at zero, the mean time from \(V=v\) is
\(\sum_{k=1}^v(\kappa k)^{-1}=\kappa^{-1}\log v+O(1)\).

\subsection{2. The unconditional B
episode}

Let \(t\in L_B\) be any actual target. Strong connectivity gives a
simple directed path

\[
       t=y_0\longrightarrow y_1\longrightarrow\cdots
                         \longrightarrow y_m=c.       \tag{2.1}
\]

Follow its designated labels with every physical clock retained, stop on
the first competitor, and, on reaching \(c\), take one final ordinary
all-clock jump. There is no structural-exit test. On designated success
the pre-final population is

\[
                         z=x-t+c\ge c.                \tag{2.2}
\]

Equation (1.2) and divergence of \(X\) make \(q\) enabled at \(z\).
Since the inactive coordinates differ from their start values by a
bounded amount,

\[
 p_c(z)\le \frac{K_c(z)_c}{K_q(z)_q}
       \le \frac{C(1+U+V)^2}{X}.                    \tag{2.3}
\]

For one ordinary marked jump put \(D(x,t)=\mathbb E_{x,t}\Delta F\).
Exact factorial cancellation gives

\[
 D(x,t)=\log p_t-\sum_y p_y\log p_y-\log K_t+\sum_y p_y\log K_y
       \le\log p_t+C.                               \tag{2.4}
\]

If \(a_i\) is the probability of the designated label at stage \(i\),
the literal all-clock Bellman recursion is

\[
 J_m=D_m,\qquad J_i=D_i+a_iJ_{i+1}.                  \tag{2.5}
\]

The first source probability tending to zero gives the negative term;
all earlier designated probabilities stay bounded below, and the later
positive tail is multiplied by the rare designated probability.
Equations (2.3)--(2.5) give \(J_0\to-\infty\).

This already proves the qualitative unconditional B episode whenever
\(\log(1+U)+\log(1+V)=o(\log X)\), including the immediate \(q=X\) case
and the B/B case with two subpower inactive coordinates.

We need one quantitative refinement only for a B episode
\textbf{appended after the prefix absorption of Section 3}. Before
absorption \(U=0\), so the absorbing target leaves \(U\) bounded by two;
\(V\), denoted \(v\), is the only possibly large inactive coordinate.
Along (2.1), every enabled source propensity is, up to a fixed factor,

\[
                         X^a(1+v)^b,
 \quad a\in\{0,1\},\quad b\in\{0,1,2\},\quad a+b\le2. \tag{2.6}
\]

Here is the complete first-rare product argument. Expand (2.5):

\[
 J_0=D_0+a_0D_1+\cdots+
             (a_0\cdots a_{m-1})D_m.                \tag{2.6a}
\]

After a subsequence, either \(v\) is bounded or \(v\to\infty\), and
every enabled source in every one of the finitely many
bounded-displacement states has a fixed exponent pair \((a,b)\) from
(2.6). The total hazard is comparable with the lexicographically largest
enabled pair, where the \(X\)-exponent is compared first because
\(v=X^{o(1)}\). Consequently a source probability is either bounded
below, at most \(C/v\), or at most \(Cv^2/X\).

Let \(j<m\) be the first path source whose probability is not bounded
below. Then \(a_0\cdots a_{j-1}\ge b_*>0\), while (2.4) gives

\[
 D_j\le
 \begin{cases}
  -\log v+C,&p_{y_j}\le C/v,\\
  -\log X+2\log v+C,&p_{y_j}\le Cv^2/X.
 \end{cases}                                          \tag{2.6b}
\]

Every \(D_i\le C_0\), and the complete positive tail after \(j\) is
multiplied by \(a_j\le Cp_{y_j}\). Hence it is \(O(p_{y_j})\), while the
negative term in (2.6b) is multiplied only by the fixed \(b_*\). If no
such \(j\) exists, all preterminal \(a_i\)'s are bounded below and the
terminal estimate (2.3) supplies the second line of (2.6b). Since
\(\log(1+v)=o(\log X)\), this proves constants \(a,C>0\), uniform over
the finite target/path menu, such that

\[
 \begin{cases}
 J_0\le-a\log X+C,&v\le M,\\
 J_0\le-a\log(2+v)+C,&v>M
 \end{cases}                                          \tag{2.7}
\]

for a fixed sufficiently large \(M\). This is only the integer-exponent
comparison (2.6); no population ranges are enumerated. Source entropy
gives every fixed positive reward moment, and the episode has at most
ten physical jumps with uniformly bounded duration moments.

\subsection{3. Exact reduction before
access}

If the witness has a cofactor, binaryity and relabelling give

\[
                              q=X+U.                  \tag{3.1}
\]

If \(q=X\), or if \(q=X+U\) and \(U>0\), the witness is already enabled.
The actual mark is still in \(L_0\), so one does \textbf{not} start the
same-linkage path (2.1) from that mark. Instead take one ordinary
all-clock payoff jump. The marked source has rate at most
\(C(1+U+V)^2\), whereas \(q\) has rate at least \(cX\). Thus (2.4) gives

\[
             \mathbb E\Delta F\le
                -\log X+2\log(2+U+V)+C.             \tag{3.2}
\]

It remains to have \(q=X+U\) and start with \(U=0\). Stop the
\textbf{prefix} at the first of:

\begin{enumerate}
\def\labelenumi{\arabic{enumi}.}
\tightlist
\item
  a reaction of \(L_B\); or
\item
  a reaction of \(L_0\) whose target contains \(U\).
\end{enumerate}

Before that event, \(X\) and \(U\) are constant, and every enabled
source and target lies in

\[
                             E=\{0,V,2V\}.            \tag{3.3}
\]

The actual mark also lies in \(E\). The nonabsorbed marked phase is

\[
              \mathcal E=\{(v,s):v\in\mathbb N_0,\ 
                              s\in\{0,1,2\},\ v\ge s\}, \tag{3.4}
\]

where \(s\) is the \(V\)-degree of the actual target. No cap replaces
\(v\).

At a type-1 event, include the causing reaction and start Section 2 from
its actual B target. At a type-2 event, include the causing reaction and
then take the access jump in (3.2), unless the causing reaction is
itself a B reaction, in which case type 1 has priority. These are the
disjoint absorbing payoff kernels.

\subsection{4. The marked killed
resolvent}

Let \(P\) be the exact substochastic embedded kernel of nonabsorbed
transitions in (3.4), and let \(S\) be its physical absorption kernel,
including the causing reaction. For \(0<\theta<1/4\), put

\[
                         w_\theta(v,s)=((v-s)!)^\theta. \tag{4.1}
\]

On a terminal state reached with source \(y\), write \(v=V\) for its
pre-jump inactive count and define

\[
 \widehat w_\theta(y;X,v)
   =\left\{\frac{(X-y_X)!\,(v-y_V)!}{X!}\right\}^{\!\theta}. \tag{4.2}
\]

Every enabled source before absorption has \(y_U=0\). Thus (4.2) is
exactly the terminal marked factorial divided by the constant prefix
factor \((X!)^\theta\); the actual target cancels. A B source of
\(X\)-degree one supplies the favorable factor \(X^{-\theta}\).

\begin{quote}
\textbf{Lemma 4.1 (uniform relative marked-factorial resolvent).} If the
phase has no reachable closed no-absorption class, then \[
      (I-P)^{-1}S\widehat w_\theta(\,\cdot\,;X,\,\cdot\,)(v,s)
         \le C_\theta w_\theta(v,s)                 \tag{4.3}
\] for all \(X\ge1\) and \((v,s)\in\mathcal E\). Equivalently, at the
physical prefix endpoint \(\sigma\), \[
      \mathbb E_{x,t}
         \exp\{\theta(F(X_\sigma,T_\sigma)-F(x,t))\}
            \le C_\theta.                           \tag{4.4}
\]
\end{quote}

\subsubsection{Proof}

Delete zero-vector self labels. Let \(d\) be the largest \(V\)-degree of
an enabled degree-zero source occurring in either linkage. There is only
one one-species complex of degree \(d\). If it belongs only to \(L_B\),
every firing from it is absorbed and the continuation is contractive.
Otherwise it is the maximal Flat0 source \(dV\).

For a nonabsorbed jump with source \(yV\), factorial cancellation gives

\[
 \frac{w_\theta(v-y+u,u)}{w_\theta(v,s)}
       =\left(\frac{(v-y)!}{(v-s)!}\right)^\theta.    \tag{4.5}
\]

The target \(u\) cancels. When \(y<d\), divide its source rate by the
order-\((v)_d\) maximal-source rate. If \(y<s\),

\[
 \frac{K_y(v)_y}{K_d(v)_d}
 \left(\frac{(v-y)!}{(v-s)!}\right)^\theta
       \le C v^{-(d-s)-(1-\theta)(s-y)}.             \tag{4.6}
\]

If \(s\le y<d\), the exponent is \(-(d-s)+(1-\theta)(y-s)<0\). Hence
every lower-source continuation is \(O(v^{-\eta})\), unless \(s<d\), in
which case the maximal-source jump itself contracts by
\(O(v^{-\theta(d-s)})\).

Only \(s=d\) permits one weight-neutral maximal-source jump. Strong
connectivity supplies an outgoing nonself label from \(dV\). It is
either absorbed or has target \(u<d\); in the latter case the successor
mark is strictly lower, so the next maximal-source jump contracts.
Consequently, for a kernel \(Q\) retaining maximal-source nonabsorbed
jumps, there are \(M<\infty\), \(\rho<1\), and \(\eta>0\) such that

\[
       \|Q^2\|_{w_\theta,\{v>M\}}\le\rho,\qquad
       \|P-Q\|_{w_\theta,\{v>M\}}\le C M^{-\eta}.     \tag{4.7}
\]

Choose \(M\) so the second bound is smaller than \((1-\rho)/4\).

Absorbing degree-one B sources cannot spoil the estimate. If \(r\) is
their propensity divided by the propensity of the enabled marked source,
their source probability is at most \(Cr/(1+r)\), while the exact marked
factorial ratio is at most \(Cr^{-\theta}\). Their weighted terminal
contribution is at most

\[
                         C\frac{r^{1-\theta}}{1+r}\le C. \tag{4.8}
\]

The identical sourcewise estimate handles every degree-zero absorbing
label. Thus the terminal operator from \(w_\theta\) to
\(\widehat w_\theta\) has finite weighted norm.

On the finite marked core \(v\le M+4\), absence of a reachable closed
no-absorption class says that the substochastic kernel is transient. For
a fixed core phase, write the total Flat0 continuation rate as \(a\),
the total absorbing degree-zero rate as \(b\), and the total coefficient
of absorbing degree-one B rates as \(c\). Its nonabsorbed row is the
fixed Flat0 numerator divided by \(a+b+cX\), hence is entrywise
nonincreasing in \(X\). The entrywise largest core kernel occurs at
\(X=1\), and its finite Green matrix is bounded. Combining that inverse
with (4.7), or using the Schur complement at first core return, gives

\[
                         \|(I-P)^{-1}\|_{w_\theta}<\infty. \tag{4.9}
\]

Equations (4.8)--(4.9) prove (4.3), and exact factorial cancellation
turns (4.3) into (4.4). \(\square\)

For every fixed \(m\), (4.4) implies

\[
          \sup_{X,v,s}\mathbb E[((F_\sigma-F_0)^+)^m]<\infty.
                                                               \tag{4.10}
\]

This uniform relative statement is what permits a Bellman decrement
which may diverge arbitrarily slowly.

We also need the continuation maximum before absorption, without the
favorable \(X^{-\theta}\) factor at a degree-one B terminal. For
\(M>v+4\), let \(\rho_M\) be the first nonabsorbed phase with
\(V\ge M\). Stop the Green series at \(\rho_M\wedge\sigma\). The same
outer contraction (4.7) and finite-core inverse (4.9), now with boundary
payoff \(w_\theta\) at \(\rho_M\) and zero at absorption, give

\[
 \mathbb E_{v,s}\!\left[
    w_\theta(V_{\rho_M},S_{\rho_M});\rho_M<\sigma
 \right]\le C_\theta w_\theta(v,s).                 \tag{4.11}
\]

Since jumps are bounded and
\(w_\theta(M-O(1),2)\ge\exp\{cM\log(2+M)\}\), (4.11) implies

\[
 \mathbb P_{v,s}\!\left\{\sup_{r<\sigma}V_r\ge M\right\}
 \le C\exp\{\theta v\log(2+v)-cM\log(2+M)\}.         \tag{4.12}
\]

This is the required survival-before-absorption tail. It is stronger
than the terminal relative bound (4.4) in precisely the case where a
terminal B source contributes \(X^{-\theta}\).

\subsection{5. Physical duration}

\begin{quote}
\textbf{Lemma 5.1 (sharp logarithmic duration).} Under Lemma 4.1, for
every fixed \(m\ge1\), \[
                   \mathbb E_{v,s}\sigma^m
      \le C_m\{1+\log(2+v)\}^m.                    \tag{5.1}
\]
\end{quote}

\subsubsection{Proof}

Put \(L(v)=1+\log(v+e)\), and give an absorbed state value zero. If the
largest source degree is one, its nonabsorbed strict jump lowers \(v\)
by one and has rate at least \(cv\); every upward jump has constant
rate. If the largest degree is two, a maximal-source nonabsorbed strict
jump lowers \(v\) by one or two and has rate at least \(cv^2\), whereas
every upward jump has source degree at most one. A maximal-source
absorbing jump adds a negative killed term. Bounded-jump Taylor
estimates give, outside a fixed core,

\[
                        {\cal K}L(v)\le-c,             \tag{5.2}
\]

where \({\cal K}\) is the physical killed generator. In degree zero,
either absorption has a fixed positive hazard or the phase is closed.
Finite-core killed transience supplies (5.2) globally after adding a
bounded finite-core corrector.

Dynkin gives (5.1) for \(m=1\). Applying the calculation to \(L^j\) and
inducting with the strong Markov property gives
\({\cal K}L^j\le-c_jL^{j-1}+C_j1_{\{v\le M\}}\), hence all integer
moments. Other fixed moments follow by domination. \(\square\)

\subsection{6. The closed no-kill
alternative}

Suppose the nonabsorbed phase has a reachable closed communicating
class. There every physical reaction preserves \(X\), preserves \(U=0\),
and remains in (3.3). It is a closed physical subset of \(\Gamma\).
Irreducibility forces it to be the whole class.

The remaining coordinate is a one-species binary mass-action chain. If
its largest nontrivial source degree is one, strong connectivity
supplies a strict \(V\to0\) edge of linear rate and every upward edge
has constant rate. If it is two, strong connectivity supplies a strict
edge from \(2V\) to a lower complex of quadratic rate, while every
upward edge has source degree at most one. Therefore, for
\(0<\theta<1/2\),

\[
 \Psi_\theta(v)=\exp\{\theta v\log(v+e)\}
\]

satisfies

\[
        {\cal L}\Psi_\theta(v)
          \le C-c(1+v)^d\Psi_\theta(v),
          \qquad d\in\{1,2\},                       \tag{6.1}
\]

outside a finite set. Degree zero is finite. Ordinary Foster theory
proves nonexplosion and positive recurrence. This is alternative 1 of
Theorem 1.1, with no episode handoff.

\subsection{7. Completion and the physical-time
ledger}

Assume the transient alternative. At type 1, append Section 2 from the
actual B target. At type 2, append the ordinary access jump (3.2). Call
the resulting stop \(\tau\). Every jump is counted once: the
absorption-causing jump ends the killed prefix, and the appended rule
starts only at its actual endpoint.

Let \(r_0=V_0-(t_0)_V\), and let \(r_\sigma\) be the residual inactive
population at the prefix endpoint. On \(r_\sigma<r_0/2\), exact
factorial telescoping gives

\[
 F_\sigma-F_0\le-c r_0\log(2+r_0)+C.                \tag{7.1}
\]

The appended B/access rule has expected reward bounded above by a fixed
constant, because every one-jump \(D\le C\) and its Bellman recursion
has finite depth. Thus the full completed reward obeys the same estimate
on this event, after changing \(C\).

On \[
             r_0/2\le r_\sigma\le2r_0+4,             \tag{7.2a}
\] (2.7) or (3.2) gives conditional appended reward at most

\[
                         -a\log(2+r_0)+C              \tag{7.2}
\]

when \(r_0\to\infty\). On \(r_\sigma>2r_0+4\), the universal upper bound
\(D\le C\), the finite Bellman recursion, and source entropy bound the
appended expected reward above by a constant. Moreover (4.12), with
\(M=2r_0+O(1)\), gives

\[
 \mathbb P\{r_\sigma>2r_0+4\}
 \le C\exp\{-c r_0\log(2+r_0)\}.                    \tag{7.2b}
\]

The positive prefix overshoot is uniformly integrable by (4.10). The
low, middle, and high events form a partition; (7.2b) makes the high
probability vanish. Therefore

\[
              \mathbb E[F(X_\tau,T_\tau)-F(x,t)]
                 \le-a_0\log(2+V_0)+C.              \tag{7.3}
\]

If \(V_0\) is bounded, (4.12) lets us choose a fixed \(M\) such that the
prefix endpoint lies below \(M\) with probability at least \(3/4\),
uniformly in \(n\). On that event (2.7) or (3.2) contributes
\(-a\log X+C\). On the complement, the positive part of the complete
reward is uniformly integrable by (4.10) and source entropy. Thus

\[
              \mathbb E[F(X_\tau,T_\tau)-F(x,t)]
                 \le-a_1\log X+C.                   \tag{7.4}
\]

Equations (5.1), (7.3), and (7.4) permit one fixed sufficiently small
\(\eta>0\) and give (1.5). Equation (4.10), the finite-jump
source-entropy bound, and Hölder give (1.6). Every endpoint is physical
and the physical duration is integrable.

The menu is finite: choose the witness, the simple B path, and tie
ordering from the finite support. If uniform constants failed, extract a
one-active subsequence with fixed actual mark, witness, path, and
bounded reaction offsets. Sections 2--7 contradict the failure. This
proves Theorem 1.1.

\subsection{8. Exact scope and the remaining global
interface}

Theorem 1.1 is the cap-free one-active B/F0 contract. The qualitative
part of Section 2 also proves the unconditional B/B contract with two
subpower inactive coordinates: every actual target belongs to a B
linkage, and (2.3) tends to zero. The two-active Q/U/C bridge and the
same cap-free path argument prove the unconditional AA contract.

This note does not alone assert recurrence of every 6,654 failure pair.
A final pair theorem must use one proper potential on complementary
passing and all-active sequences, or prove an episode-level population
boundary estimate lifting these marked episodes to that common
potential. No terminal-chart exit, descriptor SCC, or unweighted
drift-or-exit gluing is claimed here.
